
\subsubsection{Interpretation}

The results indicate that pre-trained projection weights in 7B-scale Transformers do not exhibit low-rank structure at the tested variance threshold. The $50\times$ gap between LoRA's typical adaptation rank (r ~ 32) and measured weight rank (r\_eff ~ 1600) suggests these are independent properties.

One interpretation: pre-training on diverse corpora may require high-dimensional weight structure to support varied downstream tasks. Task-specific fine-tuning identifies low-dimensional subspaces relevant to particular applications. This would explain both LoRA's success (low-rank updates for specific tasks) and our findings (high-rank weights for general capabilities).

The near full-rank structure (r\_eff ~ 40\% of dimension) is consistent with distributed representation theories where semantic information spreads across dimensions rather than concentrating in fewer principal components.

\subsubsection{Implications}

\textbf{For post-hoc conversion techniques:} Methods assuming bounded-state compression (r\_eff < 256) face an empirical constraint at 7B scale. The measured ranks are not moderately above threshold---they approach near full-rank.

\textbf{For parameter-efficient fine-tuning:} Results clarify that LoRA exploits low-dimensional structure of task-specific variations, not compression of already-low-rank weights. Rank selection strategies should focus on task properties rather than assumed weight properties.

\textbf{For compression research:} Techniques assuming weight-level low-rank structure should validate assumptions empirically before implementation at this scale.

\subsubsection{Limitations and Future Directions}

\textbf{Cross-scale analysis:} Extend to GPT-2 (117M), Pythia-1B, LLaMA-13B, LLaMA-70B to determine whether rank scales with model dimension or plateaus.

\textbf{Runtime attention analysis:} Measure effective rank of runtime attention matrices (QK^T) during inference. Runtime patterns may exhibit low-rank structure even if projection weights are full-rank.

\textbf{Architecture generalization:} Test Vision Transformers, encoder-decoder models, and multilingual variants to identify whether findings generalize beyond decoder-only language models.

\textbf{Mechanistic investigation:} Investigate why task-specific updates occupy dramatically lower-dimensional subspaces than pre-trained representations.

\subsubsection{Broader Context}

This measurement study establishes empirical ground truth for one structural property. The findings do not preclude all compression methods---only those explicitly assuming low-rank projection weights at this scale. Methods based on sparsity, quantization, or behavioral distillation operate on different assumptions.

The work demonstrates value in directly measuring assumptions rather than inferring properties from method success. LoRA's effectiveness with low-rank updates did not imply low-rank weights, though the inference was tempting. Direct measurement resolved the question.
